\section{Scope limits and non-claims}
\label{sec:scope-limits}

This section delineates explicitly what this paper does \emph{not}
claim. These limits are not accidental omissions, but necessary
structural consequences of the chosen formal scope and of the
constraints imposed by the Ontology of Continua (OC Core) and
ETS.K\_EA.v1.1.

\subsection{No inference or estimation}
\label{subsec:no-inference}

This work does not introduce, assume, or imply any estimator,
reconstruction method, or inference procedure. In particular:
\begin{itemize}
\item no mapping from observable signatures to internal enterprise
      states is constructed or postulated;
\item no decision procedure, algorithm, test, or diagnostic workflow is
      proposed;
\item no notion of convergence, sample size, learning, information
      accumulation, or observer improvement is employed.
\end{itemize}

Decidability is treated purely as a logical property of predicates under
observational equivalence. It is not framed as a capability of any
method, observer, organization, or analytical technique.

\subsection{No probabilistic or statistical guarantees}
\label{subsec:no-probability}

No probabilistic or statistical framework is used or implied. In
particular:
\begin{itemize}
\item no noise models, likelihoods, priors, posterior beliefs, or
      stochastic dynamics are assumed;
\item no confidence levels, error bounds, convergence rates, or
      asymptotic guarantees are stated;
\item no probabilistic relaxation, approximation, or high-probability
      variant of undecidability is admitted.
\end{itemize}

All results are absolute within the formalism: a predicate is either
invariant across an observational equivalence class or it is not.
No graded, approximate, probabilistic, or asymptotic notion of
decidability is admissible under the constraints of this paper.

\subsection{No intervention, recovery, or control semantics}
\label{subsec:no-control}

The paper makes no claims about how enterprises should be influenced,
stabilized, recovered, or controlled. Specifically:
\begin{itemize}
\item no intervention operator is defined or assumed;
\item no recovery path from INERTIA, COLLAPSE, or STOP is discussed;
\item no control logic, policy design, feedback mechanism, or steering
      strategy is introduced.
\end{itemize}

The results concern limits of state discrimination only. They do not
address how enterprise trajectories might be altered, optimized, or
managed after a state has been classified.

\subsection{No governance, optimization, or KPI framing}
\label{subsec:no-governance}

This work does not engage with governance frameworks, optimization
criteria, or performance management narratives. In particular:
\begin{itemize}
\item no KPIs are defined beyond ETS-admissible observability constructs;
\item no objective functions, cost functions, utility measures, or
      optimality criteria are introduced;
\item no managerial, organizational, normative, or prescriptive
      interpretation is attached to decidability results.
\end{itemize}

Decidability is not equated with usefulness, desirability, or managerial
actionability. A predicate may be decidable yet operationally trivial,
or undecidable yet strategically relevant.

\subsection{Logical, not practical, limits}
\label{subsec:logical-limits}

The limits established in this paper are logical and structural, not
practical, empirical, or technological. Undecidability is not
attributed to insufficient data, inadequate tooling, imperfect models,
or methodological immaturity.

It follows directly from the structure of the OC/ETS formalism under
incomplete observability. Accordingly, no extension of data collection,
analytics, modeling, or computation \emph{within the same formal
framework} can overcome the undecidability results presented here.
Any such extension would constitute a change of formal system and
therefore fall outside the scope of the present work.
